% ============================================================
\chapter{Referência rápida}
\label{ap:referencia}
% ============================================================

\section*{Python básico (Aulas 1 e 2)}

\begin{multicols}{2}
\begin{codigo}
# entrada e saída
n = int(input("n = "))
x = float(input("x = "))
print("resultado:", x, n)
print(f"{x:.3f}  {x:e}")

# operações
a + b, a - b, a * b, a / b
a // b      # quociente inteiro
a % b       # resto
a ** b      # potência
abs(a), round(a, 2)

# comparação e lógica
==  !=  <  <=  >  >=
and   or   not

# decisão
if cond:
    ...
elif outra:
    ...
else:
    ...
\end{codigo}

\columnbreak

\begin{codigo}
# repetição
for i in range(1, n + 1):
    ...
while cond:
    ...
break, continue

# listas
L = [3, 1, 2]
L[0], L[-1], L[1:3]
L.append(5); L.insert(0, 9)
len(L), sum(L), max(L)
sorted(L), L.copy()
[k**2 for k in range(10)]

# tuplas e dicionários
P = (1, 2);  x, y = P
D = {"a": 1};  D["a"]

# funções
def f(x, a=1):
    return a * x**2
g = lambda x: x + 1
\end{codigo}
\end{multicols}

\section*{Módulo \texttt{math} e NumPy (Aulas 3 a 5)}

\begin{multicols}{2}
\begin{codigo}
import math
math.sqrt(x), math.exp(x)
math.log(x), math.log(x, b)
math.sin(x), math.radians(g)
math.pi, math.e
math.floor(x), math.ceil(x)
math.factorial(n)
math.comb(n, k), math.gcd(a, b)
math.isqrt(n)
\end{codigo}

\columnbreak

\begin{codigo}
import numpy as np
v = np.array([1, 2, 3])
np.arange(a, b, h)
np.linspace(a, b, n)
v * 2, v + v, v @ v
np.sin(v), np.exp(v)
v[v > 0], np.sum(v), v.mean()
np.cumsum(v), np.argmax(v)
\end{codigo}
\end{multicols}

\begin{codigo}
A = np.array([[1, 2], [3, 4]])
A.shape, A[0, 1], A[:, 0], A.T, A @ B
np.linalg.det(A), np.linalg.inv(A), np.linalg.solve(A, b), np.linalg.eig(A)
np.roots([1, -3, 2]), np.polyval(p, x), np.polyfit(x, y, grau)
\end{codigo}

\section*{Matplotlib (Aula 6)}

\begin{codigo}
import matplotlib.pyplot as plt
x = np.linspace(a, b, 400)
plt.plot(x, f(x), "r--", label="f")      # cor/estilo, legenda
plt.title("..."); plt.xlabel("x"); plt.ylabel("y")
plt.xlim(a, b); plt.ylim(c, d); plt.grid(True); plt.legend(); plt.axis("equal")
plt.savefig("figura.png", dpi=200); plt.show()

fig, ax = plt.subplots(1, 2, figsize=(10, 4))   # vários gráficos: ax[0].plot(...)
X, Y = np.meshgrid(x, y); plt.contour(X, Y, Z, levels=[0])   # curvas de nível
ax = plt.figure().add_subplot(projection="3d"); ax.plot_surface(X, Y, Z)
\end{codigo}

\section*{SymPy (Aulas 7 a 10)}

\begin{codigo}
import sympy as sp
x, y = sp.symbols('x y')                       # declarar símbolos
sp.Rational(1, 3), sp.sqrt(2), sp.pi, sp.E, sp.I, sp.oo
e.evalf(20), sp.N(e, 20), e.subs(x, 2), e.subs({x: 1, y: 2})
sp.expand(e), sp.factor(e), sp.simplify(e), sp.apart(e), sp.collect(e, x)
sp.solve(e, x), sp.solve(sp.Eq(a, b), x), sp.solveset(e, x, sp.S.Reals)
sp.linsolve([e1, e2], [x, y]), sp.nsolve(e, x, chute)
sp.limit(f, x, a, '+-'), sp.limit(f, x, sp.oo), sp.diff(f, x, n)
sp.integrate(f, x), sp.integrate(f, (x, a, b)), sp.summation(f, (k, 1, n))
sp.series(f, x, x0, n).removeO()
M = sp.Matrix([[1, 2], [3, 4]]); M.det(), M.inv(), M.rref(), M.eigenvects()
y = sp.Function('y'); sp.dsolve(sp.Eq(y(x).diff(x), y(x)), y(x), ics={y(0): 1})
sp.plot(f, (x, a, b)); sp.lambdify(x, f, "numpy")
\end{codigo}

\section*{Mensagens de erro mais comuns}

\begin{center}\small
\begin{tabularx}{\linewidth}{@{}l>{\raggedright\arraybackslash}X@{}}
\toprule
\textbf{Erro} & \textbf{Causa provável}\\
\midrule
\py{SyntaxError} & faltam dois-pontos, parênteses ou aspas; uso de \py{^}; vírgula decimal\\
\py{IndentationError} & recuo faltando ou inconsistente após \py{if}, \py{for}, \py{def}...\\
\py{NameError: name 'x' is not defined} & variável/símbolo não criado, célula anterior não executada, erro de digitação\\
\py{TypeError} & operação entre tipos incompatíveis (texto + número; \py{math.sin} de um \emph{array})\\
\py{ValueError} & valor inválido: \py{int("3.5")}, \py{math.sqrt(-1)}, limite que não existe\\
\py{IndexError} & índice fora da lista (lembre que começa em 0)\\
\py{ZeroDivisionError} & divisão por zero\\
\py{ModuleNotFoundError} & biblioteca não instalada: \texttt{\%pip install nome}\\
\bottomrule
\end{tabularx}
\end{center}
